\begin{lemma}[Hereditary countability of the front recursion]
For every $s\in T(F)$, the recursively defined value
$\widetilde f(s)$ belongs to the hereditary countable part of the full
hierarchy:
\[
  \mathsf{HereditarilyCountable}(\widetilde f(s)).
\]
\end{lemma}
\begin{proof}
Proceed by well-founded induction on the extension-oriented relation of
\noderef{FrontTreeWellFounded}.  If $s\in F$, the front equation
\noderef{TildeFFrontEquation} identifies the value with an atom, and atoms
are hereditarily countable by \noderef{HereditarilyCountableAtom}.
If $s\notin F$, the non-front equation
\noderef{TildeFNonfrontEquation} identifies the value with a full node
indexed by the nonempty type $E_s$ of ordered one-point extensions.  For
each $n\in E_s$, the recursive induction hypothesis gives hereditary
countability of $\widetilde f(s\cup\{n\})$.  Thus the induction hypotheses
form a family
\[
 g:E_s\longrightarrow\mathsf{HerCtblPower}(Q),\qquad
 g(n)=\bigl(\widetilde f(s\cup\{n\}),
             \text{its induction witness}\bigr),
\]
using the subtype definition \noderef{HerCtblPower}.  The extension lemma
\noderef{FrontTreeImmediateExtensions} gives that $E_s$ is nonempty and
countable; the same is true of its universe lift.  Transport $g$ along the
lift's down projection and apply the countable hereditary-node closure theorem
\noderef{HerCtblNodeClosure}.  It proves hereditary countability of exactly
the node in the non-front recursion equation.  These are the two recursion
cases, so the assertion holds for all tree nodes.
\end{proof}
